\section{Integral allocations and a near-core construction}\label{app:allocation}
The request-cover and integral-allocation principles of
Section~\ref{sec:prelim} apply throughout this appendix.

\begin{lemma}[Exact boxed-triangle allocation]\label{lem:triangle}
Let $M_1,M_2,M_3$ be nonnegative integers and let $L_1,L_2,L_3$ be
integers satisfying $L_i\le M_j+M_l$ whenever $\{i,j,l\}=\{1,2,3\}$.
The minimum of $a_1+a_2+a_3$ over integral vectors satisfying
\[
 0\le a_i\le M_i,\qquad a_j+a_l\ge L_i
\]
is the ceiling of
\begin{equation}\label{eq:triangleexact}
 \max\left\{0,L_1,L_2,L_3,
 L_1+L_2-M_3,L_1+L_3-M_2,L_2+L_3-M_1,
 \frac{L_1+L_2+L_3}{2}\right\}.
\end{equation}
\end{lemma}
\begin{proof}
For a proposed nonnegative integer total $t$, the conditions are equivalent to
\[
 0\le a_i\le\min(M_i,t-L_i),\qquad \sum_i a_i=t.
\]
Integral coordinates exist if and only if all three upper bounds are
nonnegative and their sum is at least $t$. Expand the sum of the three
minima as the minimum of its eight affine expressions. The resulting
conditions are
\begin{gather*}
 t\ge0,L_1,L_2,L_3,\qquad
 t\ge L_i+L_j-M_l\quad(\{i,j,l\}=\{1,2,3\}),\\
 2t\ge L_1+L_2+L_3,\qquad t\le M_1+M_2+M_3,\qquad
 L_i\le M_j+M_l.
\end{gather*}
The last three host inequalities hold by assumption. Every lower bound
in \eqref{eq:triangleexact} is at most $M_1+M_2+M_3$, so the ceiling of
the maximum also satisfies the upper bound on $t$. Coordinates of this
total can be filled successively within their integral upper bounds.
The same argument without ceilings gives the real minimum.
\end{proof}

\begin{lemma}[Exact symmetric static allocation]\label{lem:symmetric}
Let a good triple have pair sizes $m\ge y\ge z\ge0$, and put $S=m+y+z$.
The four-request construction below can be implemented at integer
budget $T$ if and only if
\begin{equation}\label{eq:symexact}
 T\ge\max\left\{k+m-y-z,\frac{k+m}{2},
             \frac{2k+m+y-z}{3},\frac{3k+S}{5}\right\}.
\end{equation}
Whenever it can be implemented, it closes the triple.
\end{lemma}
\begin{proof}
Split the pair cells $M,Y,Z$ into selected parts of sizes $a,b,c$ and
their complements. With the private cells denoted by $P_E,P_F,P_G$,
use four requests
\begin{align*}
 R_0&=M_1\cup Y_1\cup Z_1,\\
 R_M&=M\cup P_G\cup Y_2\cup Z_2,\\
 R_Y&=Y\cup P_F\cup M_2\cup Z_2,\\
 R_Z&=Z\cup P_E\cup M_2\cup Y_2.
\end{align*}
Their sizes are $a+b+c$, $k+m-b-c$, $k+y-a-c$, and $k+z-a-b$.
Apply Lemma~\ref{lem:triangle} with hosts $(m,y,z)$ and requirements
\[
 L_1=k+m-T,\qquad L_2=k+y-T,\qquad L_3=k+z-T.
\]
The host conditions reduce to $T\ge k+m-y-z$. The individual
requirements $L_i\le T$ reduce to $2T\ge k+m$; the two-$L$ requirements
reduce to $3T\ge2k+m+y-z$; and the total requirement is
$5T\ge3k+S$. This proves the exact criterion.

Every pair meeting the good triple is contained in a request. A pair
cell and its opposite private part lie in the corresponding request.
For two different pair cells, if both points lie in selected parts they
belong to $R_0$; otherwise one of the last three requests contains both.
Thus four responses close the triple.

The optimality assertion concerns this specified construction only,
not all possible closing constructions.
\end{proof}

\begin{lemma}[Exact static Hall allocation]\label{lem:hallstatic}
For a good triple with pair sizes $a,b,c$, the four-request allocation
below has an integral implementation at budget $T$ if and only if
\begin{gather}
 T-a\ge0,\quad T-a-b\ge0,\quad T-a-c\ge0,\quad T-b-c\ge0,
 \label{eq:hallbase}\\
 k+2c\le2T,\quad k+2b\le2T,\quad k+3a\le3T,\quad
 2k+b+c\le3T,\label{eq:hallseven}\\
 2k+2a+c\le4T,\quad2k+2a+b\le4T,\quad3k+a\le4T.\nonumber
\end{gather}
Whenever it has such an implementation, it closes the triple.
\end{lemma}
\begin{proof}
Label the pair cells $X=E\cap F$, $Y=E\cap G$, $Z=F\cap G$,
of sizes $a,b,c$, by request-index sets $\{1,2,3\}$, $\{2,4\}$,
and $\{3,4\}$, respectively. A label specifies all requests containing
its point. Assign the private points to singleton labels according to
\[
 P_E:\{3\}\text{ or }\{4\},\qquad
 P_F:\{2\}\text{ or }\{4\},\qquad
 P_G:\{1\},\{2\}\text{ or }\{3\}.
\]
The four remaining capacities are exactly those in
\eqref{eq:hallbase}. The seven Hall inequalities for the three source
sets $P_E,P_F,P_G$, of sizes $k-a-b,k-a-c,k-b-c$, are precisely
\eqref{eq:hallseven}. Thus the integral-allocation principle proves the
claimed criterion.

Every candidate pair has intersecting labels. The three pair-cell
labels intersect pairwise, and every allowed label of a private cell
meets the label of its opposite pair cell. Consequently every candidate
pair is contained in some request. The request-cover principle gives
closure.

\end{proof}


\begin{lemma}[Near-core allocation]\label{lem:nearcore}
Suppose every pair intersection in the family is at most $m$ or greater
than $k/2$, where $m\le k/2$. Let $E,F,G$ be a good triple with pair
sizes $x=|E\cap F|$, $y=|E\cap G|$, $z=|F\cap G|$, and put
$\Delta=(x+y+z-T)_+$. The triple closes at integer budget $T$ if
$y+z\le T$ and
\begin{align}
 m+z&\le T,& k+y-z&\le T,\label{eq:ncfirst}\\
 2k+x-2z&\le2T,&2k+x+m-z&\le3T,\nonumber\\
 x+m&\le T,&&\nonumber\\
 k-x+y+\Delta&\le T,&2k-2x+z+\Delta&\le2T,\label{eq:nclast}\\
 2k-z+\Delta&\le2T,&3k-x-y+\Delta&\le3T.\nonumber
\end{align}
If $x+y+z\le T$, the last four bounds may instead be replaced by
\begin{equation}\label{eq:emptycore}
 k-x+y\le T,\qquad 2k-2x+z\le2T.
\end{equation}
\end{lemma}
\begin{proof}
Put $X=E\cap F$, $Y=E\cap G$, $Z=F\cap G$. Request an edge $H$
avoiding $Y\cup Z$ and a subset of $X$ of size
$\min(x,T-y-z)$. The only possibly nonempty triple cell is
$P=E\cap F\cap H$, of size $p\le\Delta$.
Define the disjoint cells
\begin{align*}
 X'&=X\setminus P, & A&=H\cap\bigl(E\setminus(F\cup G)\bigr),\\
 B&=H\cap\bigl(F\setminus(E\cup G)\bigr), &C&=H\cap G,\\
 W&=G\setminus(E\cup F\cup H).
\end{align*}
Write $a,b,c,w$ for the latter four sizes. Then
\begin{gather*}
 a\le k-x-y,\quad b\le k-x-z,\quad c\le k-y-z,\\
 p+a+b+c\le k,\qquad c+w=k-y-z.
\end{gather*}
The complete candidate products for pairs meeting $E,F,G,H$ are
\[
 P\times(Y\cup Z\cup C\cup W),\qquad
 X'\times C,\qquad Y\times B,\qquad Z\times A.
\]
A point in three edges lies in $P$; otherwise a covering pair must lie
in complementary pair cells. This proves completeness of the list.

\paragraph{Case 1: $p+a\le m$.}
Put $P$ in all three remaining requests, $X'$ in requests $1,3$,
$Y,B$ in request $1$, and $Z,A$ in request $2$. Partition $C$ between
requests $1,3$, and then partition $W$ among all three. The initial
loads are
\[
 L_1=x+y+b,\qquad L_2=p+z+a,\qquad L_3=x.
\]
They fit by $L_1\le k+y-z$, $L_2\le m+z$, and $L_3\le x+m$.
The first partition fits because
\[
 L_1+L_3+c=2x+y+b+c\le2k+x-2z\le2T.
\]
After that partition, the total load including $W$ is
\[
 L_1+L_2+L_3+c+w=2x+k+p+a+b\le2k+x+m-z\le3T.
\]
All capacities are integral, so the partitions exist. Each candidate
product is covered: $P$ lies in every request; the two possible labels
of $C$ meet the label of $X'$; and $Y,B$ and $Z,A$ share requests.

\paragraph{Case 2: $p+a>m$.}
The global dichotomy gives $p+a>k/2$. The traces $E\cap H$ and
$G\cap H$ are disjoint, so $c<k/2$ and hence $c\le m$.
Put $P$ in requests $1,2$, $X',C$ in request $2$, $Y,B$ in request
$1$, and $Z$ in requests $1,3$. Partition $A$ between requests $1,3$
and $W$ between requests $1,2$. The initial loads are
\[
 L_1=p+y+z+b,\qquad L_2=x+c,\qquad L_3=z.
\]
They fit by $L_1\le\Delta+k-x+y$, $L_2\le x+m$, and $L_3\le m+z$.
The three Hall conditions for the two flexible cells are
\[
 a\le2T-L_1-L_3,\qquad w\le2T-L_1-L_2,\qquad
 a+w\le3T-L_1-L_2-L_3.
\]
They follow from
\begin{align*}
 L_1+L_3+a&\le\Delta+2k-2x+z\le2T,\\
 L_1+L_2+w&=p+x+k+b\le\Delta+2k-z\le2T,\\
 L_1+L_2+L_3+a+w&=p+x+k+z+a+b
                  \le\Delta+3k-x-y\le3T.
\end{align*}
Thus an integral allocation exists. The labels of $P$ meet all allowed
labels of $Y,Z,C,W$; $X',C$ share request $2$; $Y,B$ share request $1$;
and both labels of $A$ meet the label $\{1,3\}$ of $Z$. Every candidate
pair is therefore missed by a response.

Finally, if $x+y+z\le T$, the first request omits the whole pair-cell
union, so $P=\varnothing$. The cell $W$ has no candidate-pair neighbour
and can be left out of every remaining request. In Case 2 only the
allocation of $A$ between requests $1,3$ is needed. Its conditions are
exactly the base-capacity bounds and \eqref{eq:emptycore}. Case 1 already
works even if $W$ is included. This proves the empty-core variant.
\end{proof}
